% homework/hw09/hw09.tex
% Year 7 - Homework Packet 9: Lighter on the Moon.
%
% Covers Mathematics 4.1-4.4 (compare, add, subtract, multiply and divide
% decimals), 4.5 (making calculations easier) and the Unit 4 review (x and / by
% a 2-digit number, money to 2 d.p., conjecture and the mean). Section F takes
% Science 3.1 (gravity, weight and mass) and uses it ONLY as a place to do
% decimals -- at the teacher's request, there is no science question on this
% packet that is not answered with a decimal.
%
% THE SPINE IS THAT NOTHING IN UNIT 4 IS NEW ARITHMETIC. 3.2 x 7 is 32 x 7 with
% the point put back; 8.24 / 4 is short division with the point straight up.
% The only new question is where the point goes, and the estimate answers it.
% Every section asks that question in a different way.
%
% Section F is why this packet exists. F3 is the falling-ball lab done with a
% stopwatch: four drop times each for a heavy ball and a light ball. The
% fastest SINGLE time belongs to the light ball, so a student who reads one
% number gets the science wrong. Both MEANS are 0.45 s. The decimals answer
% "do heavy things fall faster?", which is the question 3.1 opened on, and the
% spread of the times is the 0.2 s reaction time from slide 43.
%
% Nothing untaught:
%   * Strengths of gravity are the 3.1 deck's own (widgets.jsx, Mr Bowen in
%     Space): Moon 1.6, Mars 3.7, Uranus 8.7, Venus 8.9, Earth 10, Neptune 11,
%     Jupiter 25. Mercury (3.7) and Saturn (10) are left out so F4 has ONE
%     answer.
%   * No decimal x decimal anywhere: every weight is a whole-number mass times
%     a decimal strength, which is 4.3 and the review. "A quarter of 0.5" in E4
%     is 0.5 / 4, which is 4.4.
%   * The deck already did Mr Bowen's 80 kg on every planet, so F2 uses a 15 kg
%     backpack, and F4 runs the deck's question backwards.
%   * No number is lifted from the workbook or from the decks' practice panels.
%
% House style is shared: see homework/hw-style.tex. Method: the new-homework
% skill. Source is pure ASCII on purpose.
%
% Build:  pwsh .claude/skills/new-homework/scripts/build-hw.ps1 hw09

\documentclass[12pt,a4paper]{article}

% -- Answer key switch -------------------------------------------------------
\newif\ifkey
\ifdefined\TEACHER\keytrue\else
  \keyfalse        % \keyfalse = student copy   |   \keytrue = teacher copy
\fi

% -- Language switch ---------------------------------------------------------
% The French copy (hw09-fr.pdf) is this file built with \FRENCH defined. See
% hw09x.tex, which introduced the scheme, and the README.
\newif\iffr
\ifdefined\FRENCH\frtrue\fi

% -- House style -------------------------------------------------------------
% homework/hw-style.tex
% Shared house style for every Year 7 homework packet. \input this from the
% packet file, right after \documentclass and the answer-key switch.
%
% Do not fork this file per packet. If a packet needs a different look, the
% question is whether every packet should have it.
%
% The choices here are deliberate and were arrived at by printing and looking:
%   * No solid dark banners. Every header is a light tint with a coloured spine
%     on the left, so colour reads clearly without flooding a school printer.
%   * No marks and no mark totals. These are practice, not exam papers.
%   * Lato at 12pt with open leading, plus mathastext so inline maths is Lato
%     too. Without mathastext every "-6 + 4" is Computer Modern serif and the
%     sheet looks half-typeset.
%   * Writing lines are spaced for Year 7 handwriting, not adult margin notes.
%   * Pages fill. \needspace is redefined below so a heading no longer drags
%     the page break up to itself, and every heading is glued to the box under
%     it. Packets built before 2026-09-30 were laid out under the old rules and
%     will reflow if rebuilt -- re-read every page if you rebuild one.
%
% Provides: \wlines \blank \tickbox \sep \qmk-free marks (none) \hwsection
% \qhead \expr, the workspace box, and the remember / wordhelp / activity /
% yourturn coloured boxes. Palette matches the lesson decks exactly.

% -- Packages ----------------------------------------------------------------
\usepackage[T1]{fontenc}
\usepackage[utf8]{inputenc}
\usepackage[default]{lato}
\usepackage[a4paper,top=1.7cm,bottom=1.7cm,left=2.0cm,right=2.0cm,headheight=15pt]{geometry}
\usepackage{amsmath,amssymb}
\usepackage{xcolor}
\usepackage[most]{tcolorbox}
\usepackage{enumitem}
\usepackage{tikz}
\usetikzlibrary{arrows.meta}
\usepackage{array}
\usepackage{tabularx}
\usepackage{fancyhdr}
\usepackage{multicol}
\usepackage{needspace}
\usepackage{graphicx}

% Make maths use Lato too. Without this, every "-6 + 4" on the page is set in
% Computer Modern serif and the sheet looks half-typeset. Must come last.
\usepackage[italic]{mathastext}

% -- \needspace, without the hole it used to leave ----------------------------
% needspace.sty writes "\vskip 0pt plus X \penalty-100" in front of every
% heading. A page that ends anywhere else is underfull with no stretch, so TeX
% scores it b=10000; a page that ends AT a heading gets that X of stretch and a
% finite badness, and wins. So any heading in roughly the bottom half of a page
% claimed the break, and the question under it moved over whole, leaving a hole
% the size of the question. \tracingpages showed it exactly (b=4531 at the
% heading, b=10000 at every later break) while building homework/q1rev, and it
% is where the half-empty pages in HW 1-8 came from.
% This keeps only the \penalty9999, which still refuses to put a heading in the
% last X of a page. Keeping a heading WITH the box under it is \nobreak's job
% now -- see \hwsection and \qhead below.
\makeatletter
\renewcommand{\needspace}[1]{\begingroup\setlength{\dimen@}{#1}%
  \vskip\dimen@\penalty9999\vskip-\dimen@\endgroup}
\makeatother

\linespread{1.06}\selectfont

% -- The Learner's Book palette (same hex values as the lesson decks) ---------
\definecolor{cbteal}{HTML}{0087A8}
\definecolor{cbpurple}{HTML}{5C2483}
\definecolor{cborange}{HTML}{C25E12}
\definecolor{cbgreen}{HTML}{4A8B23}
\definecolor{cbred}{HTML}{C8102E}
\definecolor{cbblue}{HTML}{1A5FA8}
\definecolor{rulegray}{HTML}{9AA5AD}

% -- Answer lines and blanks -------------------------------------------------
% \wlines{n} draws n ruled writing lines, spaced wide enough for Year 7
% handwriting rather than for an adult with a fine pen.
\makeatletter
\newcommand{\wlines}[1]{%
  \par\vspace{0.75em}%
  \@tempcnta=0\relax
  \@whilenum\@tempcnta<#1\do{%
    \hbox to \linewidth{\color{rulegray}\leaders\hrule height 0.5pt\hfill}%
    \vspace{1.5em}%
    \advance\@tempcnta by 1\relax}%
  \par\vspace{-0.8em}%
}
\makeatother

% \blank[width] is an inline gap to write a short answer in.
\newcommand{\blank}[1][2.6cm]{\underline{\hspace{#1}}}

% A boxed space for showing working.
\newtcolorbox{workspace}[1][3.4cm]{%
  colback=white, colframe=rulegray, boxrule=0.5pt, arc=5pt,
  height=#1, left=6pt, right=6pt, top=4pt, bottom=4pt,
  before skip=9pt, after skip=11pt}

% An empty tick box.
\newcommand{\tickbox}{\fbox{\rule{0pt}{1.15em}\rule{1.15em}{0pt}}}

% A middot separator.
\newcommand{\sep}{\,\textperiodcentered\,}

% -- Section and question headings -------------------------------------------
% Light tint plus a fat coloured spine on the left. No solid dark banner: the
% colour reads clearly but barely costs the printer anything.
% needspace keeps a heading out of the last few lines of a page. The \nobreak
% after it keeps it with whatever comes next: a remember or word-help box that
% does not fit takes its heading over with it, instead of leaving the heading
% alone at the foot of the page. (The old needspace did that by accident, and
% left a hole every time; see above.)
\newcommand{\hwsection}[2][cbteal]{%
  \needspace{8\baselineskip}%
  \par\vspace{13pt}%
  \noindent
  \begin{tcolorbox}[enhanced, nobeforeafter, width=\textwidth,
      colback=#1!7, colframe=#1!7, arc=5pt,
      borderline west={5pt}{0pt}{#1},
      left=14pt, right=10pt, top=7pt, bottom=7pt]
    {\large\bfseries\color{#1}#2}
  \end{tcolorbox}%
  \par\nobreak\vspace{9pt}}

% \qhead[colour]{A2}{Moving along the line} -- a question heading.
\newcommand{\qhead}[3][cbteal]{%
  \needspace{6\baselineskip}%
  \par\vspace{11pt}%
  \noindent{\large\bfseries\color{#1}#2}\quad{\large\bfseries #3}\par\nobreak\vspace{4pt}}

% -- Coloured boxes ----------------------------------------------------------
% Shared look: rounded, light fill, coloured rule, coloured title text on a
% pale title strip rather than a dark one.
% NOT breakable. A breakable box that cannot fit its first lines forces its own
% page break, and \nobreak cannot stop a forced break -- so the heading above
% it was left alone at the foot of the page (tested: it happens with \nobreak
% in place). Unbreakable, the page builder moves heading and box together.
% These boxes are a few lines each; none needs to split.
\tcbset{hwbox/.style={
  enhanced, arc=6pt, boxrule=1pt,
  left=11pt, right=11pt, top=8pt, bottom=8pt,
  fonttitle=\bfseries, attach boxed title to top left={xshift=12pt, yshift=-8pt},
  boxed title style={arc=4pt, boxrule=1pt, left=7pt, right=7pt, top=2pt, bottom=2pt},
  top=13pt}}

% Orange: things to remember. Matches the deck's "Write This Down".
\newtcolorbox{remember}[1][Remember from the lesson]{hwbox,
  colback=cborange!6, colframe=cborange, coltitle=cborange,
  colbacktitle=white, title={#1}}

% Blue: English help. The barrier in this class is the wording, not the maths.
\newtcolorbox{wordhelp}[1][Word help]{hwbox,
  colback=cbblue!6, colframe=cbblue, coltitle=cbblue,
  colbacktitle=white, title={#1}}

% Purple: an activity or instruction.
\newtcolorbox{activity}[1][How to do this homework]{hwbox,
  colback=cbpurple!6, colframe=cbpurple, coltitle=cbpurple,
  colbacktitle=white, title={#1}}

% Green: your turn / a friendly nudge.
\newtcolorbox{yourturn}[1][Your turn]{hwbox,
  colback=cbgreen!7, colframe=cbgreen, coltitle=cbgreen!75!black,
  colbacktitle=white, title={#1}}

% -- Shorthands --------------------------------------------------------------
\newcommand{\dC}{\,^{\circ}\mathrm{C}}          % use inside maths: $-8\dC$
\newcommand{\key}[1]{{\color{cborange}\textbf{#1}}}
\newcommand{\eng}[1]{{\color{cbblue}\textbf{#1}}}

\setlist[enumerate]{itemsep=8pt, topsep=5pt, leftmargin=*}
\setlist[itemize]{itemsep=4pt, topsep=4pt, leftmargin=*}
\renewcommand{\arraystretch}{1.45}

% -- An expression the student is asked to work from -------------------------
\newcommand{\expr}[1]{%
  \tcbox[on line, colback=cbgreen!10, colframe=cbgreen!55, boxrule=1pt, arc=4pt,
         left=9pt, right=9pt, top=4pt, bottom=4pt]{\large\bfseries $#1$}}

% -- Running head and foot ---------------------------------------------------
% The packet overrides \fancyhead[L] with its own title.
\pagestyle{fancy}
\fancyhf{}
\fancyhead[L]{\small\color{black!55}Year 7\sep Homework}
\fancyhead[R]{\small\color{black!55}Name: \blank[3.4cm]}
\fancyfoot[C]{\small\color{black!55}Page \thepage}
\renewcommand{\headrulewidth}{0pt}
\renewcommand{\headrule}{\vspace{2pt}\hbox to\headwidth{\color{cbteal!45}\leaders\hrule height 1.2pt\hfill}}



% -- French support lines ----------------------------------------------------
% \fr{...}  a French line of its own, grey and italic, under the English.
% \frin{...} a French gloss in brackets, inline -- for table cells and labels.
% Both are EMPTY in the English copy, so the English layout is the one that was
% checked to ten pages. Key English words stay English inside the French, in
% blue, because the English is what is being learned.
% Both start with \unskip: the line break before \fr is a space, and in the
% English copy it would otherwise sit at the end of the line and move a break
% (it did, in E2(d)).
\newcommand{\fr}[1]{\unskip\iffr\par\vspace{1pt}\noindent{\small\itshape\color{black!60}#1}\par\fi}
\newcommand{\frin}[1]{\unskip\iffr\ {\small\itshape\color{black!60}(#1)}\fi}
\newcommand{\fq}[1]{\guillemotleft\,#1\,\guillemotright}

% -- Per-packet running head -------------------------------------------------
\iffr
  \fancyhead[L]{\small\color{black!55}Year 7\sep Homework 9\sep English /
    Fran\c{c}ais}
\else
  \fancyhead[L]{\small\color{black!55}Year 7\sep Homework 9}
\fi

% -- A group of writing lines that will not be split -------------------------
% See hw06.tex: \wlines is a run of separate hboxes and TeX will break in the
% middle of one, stranding a single rule at the top of the next page.
\newcommand{\qlines}[1]{\par\needspace{\numexpr#1+1\relax\baselineskip}\wlines{#1}}

% -- A box to write one sign or number in ------------------------------------
\newcommand{\nbox}{\ \fbox{\rule{0pt}{1.05em}\rule{1.0em}{0pt}}\ }

% -- An answer blank that will not cause an overfull line ---------------------
% "\hfill \blank[..]" at the end of a wrapping \item overflows the margin when
% the text happens to fill its last line. \ansline puts the blank on a line of
% its own, right-aligned, which cannot overflow.
\newcommand{\ansline}[1][4.4cm]{\par\vspace{2pt}\hspace*{\fill}\blank[#1]\par}

% -- "Calculation: ____  Answer: ____" on a line of its own -------------------
% The calculation comes FIRST on the line on purpose: it is the thing to write
% first, and the packet says so on the cover.
\newcommand{\calcans}[1][6.6cm]{\par\vspace{5pt}\noindent\hspace*{1.6em}%
  \textbf{Calculation:} \blank[#1] \hfill \textbf{Answer:} \blank[3.0cm]\par}

% -- A ragged-right X column -------------------------------------------------
\newcolumntype{L}{>{\raggedright\arraybackslash}X}

% -- A centred fixed-width column --------------------------------------------
\newcolumntype{C}[1]{>{\centering\arraybackslash}p{#1}}

% -- One calculation with a box to work it in --------------------------------
% Three to a row. The calculation sits above its own box so the two cannot be
% separated by a page break.
\newcommand{\wsq}[3][3.1cm]{%
  \begin{minipage}[t]{0.315\textwidth}
    \textbf{(#2)}\; $#3$
    \begin{workspace}[#1]\end{workspace}
  \end{minipage}}

% -- Two to a row, for the longer ones ---------------------------------------
\newcommand{\wsh}[3][3.4cm]{%
  \begin{minipage}[t]{0.485\textwidth}
    \textbf{(#2)}\; $#3$
    \begin{workspace}[#1]\end{workspace}
  \end{minipage}}

% -- A table row tall enough to write in -------------------------------------
\newcommand{\tall}{\rule[-0.6em]{0pt}{1.9em}}

\begin{document}

% ===========================================================================
% COVER BLOCK
% ===========================================================================
\thispagestyle{fancy}

\begin{tcolorbox}[enhanced, colback=cbpurple!7, colframe=cbpurple!7, arc=7pt,
                  borderline west={6pt}{0pt}{cbpurple},
                  left=18pt, right=14pt, top=10pt, bottom=10pt]
  {\color{cbpurple}\bfseries YEAR 7\sep UNIT 4\sep HOMEWORK 9}\\[5pt]
  {\color{cbpurple}\Huge\bfseries Lighter on the Moon}\\[7pt]
  {\color{black!70}Mathematics 4.1--4.5 and the Unit 4 review --- decimals}\\
  {\color{black!70}Section F --- Science 3.1 in decimals: weight, gravity and falling}
  \iffr\\[4pt]{\small\itshape\color{black!60}Plus l\'eger sur la Lune~: les
  d\'ecimaux, et la gravit\'e en d\'ecimaux}\fi
\end{tcolorbox}

\vspace{4pt}
\noindent
\begin{tabularx}{\textwidth}{@{}X X@{}}
  \textbf{Name:} \blank[5.6cm] & \textbf{Class:} \blank[5.6cm] \\
  \textbf{Date given:} \blank[4.6cm] &
  \textbf{Date due:} {\color{cbred}\bfseries Friday 23 October 2026} \\
\end{tabularx}

\vspace{2pt}
\begin{activity}
  Six sections. A to E are the decimal lessons. F is the science from 3.1, and
  \key{every answer in F is a decimal}.
  \fr{Six parties. A \`a E~: les le\c{c}ons sur les d\'ecimaux. F~: la science de la
  le\c{c}on 3.1 --- toutes les r\'eponses sont des d\'ecimaux.}
  \begin{itemize}
    \item \key{Write the calculation before you write the answer.}
          \frin{\'Ecris le calcul avant la r\'eponse.}
    \item \key{Estimate first.} The estimate tells you where the point goes.
          \frin{Estime d'abord~: l'estimation te dit o\`u va le point.}
    \item \eng{No calculator.} In F, answer in \key{full English sentences}
          where it says so.
          \frin{Pas de calculatrice. Dans F, r\'eponds par des phrases compl\`etes
          \textbf{en anglais}.}
  \end{itemize}
\end{activity}

% ===========================================================================
% SECTION A -- MATHEMATICS 4.1
% ===========================================================================
\hwsection{Section A \textperiodcentered\ Mathematics 4.1 --- Comparing and Ordering}

% -- A1 ---------------------------------------------------------------------
% The way in, and it is meant to be easy. Part (b) is the whole of 4.1 with
% the line still in front of them: once 5.05 and 5.45 are drawn, the order is
% just reading left to right.
\qhead{A1}{Put each number on the line}

\noindent The line goes from \textbf{5} to \textbf{6}. Every small line is
\textbf{0.1}. \textbf{(a)} Draw an arrow to each number and write its letter.
\fr{La droite va de 5 \`a 6. Chaque petit trait vaut 0.1. (a) Dessine une fl\`eche
vers chaque nombre et \'ecris sa lettre.}

\vspace{4pt}
\begin{center}
\begin{tikzpicture}[x=1.45cm, y=1cm, line width=0.9pt]
  \path (0,0.95) -- (10,0.95);  % room above the line for arrows and letters
  \draw[cbteal] (-0.4,0) -- (10.4,0);
  \foreach \n in {0,...,10}{\draw[cbteal] (\n,-0.15) -- (\n,0.15);}
  \draw[cbteal, line width=1.6pt] (0,-0.26) -- (0,0.26);
  \draw[cbteal, line width=1.6pt] (5,-0.22) -- (5,0.22);
  \draw[cbteal, line width=1.6pt] (10,-0.26) -- (10,0.26);
  \node[below=4pt, font=\bfseries] at (0,-0.26) {$5$};
  \node[below=4pt, font=\small] at (5,-0.22) {$5.5$};
  \node[below=4pt, font=\bfseries] at (10,-0.26) {$6$};
\end{tikzpicture}
\end{center}

\vspace{2pt}
\noindent
\begin{tabularx}{\textwidth}{|C{1.6cm}|C{1.6cm}|C{1.6cm}|C{1.6cm}|X}
  \cline{1-4}
  \textbf{A} \; $5.2$ \tall & \textbf{B} \; $5.9$ & \textbf{C} \; $5.45$ &
  \textbf{D} \; $5.05$ & \\
  \cline{1-4}
\end{tabularx}

\vspace{6pt}
\noindent \textbf{(b)} Now write the four numbers \textbf{in order of size, starting
with the smallest}.
\fr{(b) Range les quatre nombres par ordre de taille, en commen\c{c}ant par le plus
petit (\eng{starting with the smallest}).}
\ansline[9.0cm]

% -- A2 ---------------------------------------------------------------------
\qhead{A2}{Which is bigger?\frin{Lequel est le plus grand~?}}

\begin{remember}[Compare decimals]
  First look at the \key{whole-number parts}. If they are the same, \key{fill with
  zeros} so both numbers have the same number of digits after the point, then
  compare. \quad $3.3 = 3.30$, \; and $3.30 > 3.25$.
  \fr{Regarde d'abord les parties enti\`eres (\eng{whole-number parts}). Si elles sont
  pareilles, compl\`ete avec des z\'eros (\eng{fill with zeros}), puis compare.}
\end{remember}

\vspace{2pt}
% A minipage so the instruction cannot be left at the foot of a page with the
% six boxes on the next one (it was, on the first build).
\noindent\begin{minipage}{\textwidth}
Write $<$, $>$ or $=$ in each box.
\fr{\'Ecris $<$, $>$ ou $=$ dans chaque case.}

\vspace{4pt}
\begin{multicols}{3}
\begin{enumerate}[label=(\alph*), itemsep=11pt]
  \item $7.4 \;\nbox\; 7.38$
  \item $0.6 \;\nbox\; 0.59$
  \item $3.05 \;\nbox\; 3.5$
  \item $12.3 \;\nbox\; 9.99$
  \item $0.70 \;\nbox\; 0.7$
  \item $4.125 \;\nbox\; 4.13$
\end{enumerate}
\end{multicols}
\end{minipage}

\vspace{12pt}
\noindent \textbf{(g)} Mr Bowen says: \emph{``6.18 is bigger than 6.2, because 18
is bigger than 2.''}
\fr{Mr Bowen dit~: \guillemotleft\,6.18 est plus grand que 6.2, parce que 18 est plus
grand que 2.\,\guillemotright}

\vspace{3pt}
\noindent\hspace*{1.6em}(i) Write how you \textbf{say} each number in English.
\frin{Comment dit-on chaque nombre en anglais~?}

\vspace{4pt}
\noindent\hspace*{1.6em}$6.18$: \blank[5.2cm] \hfill $6.2$: \blank[5.2cm]

\vspace{6pt}
\noindent\hspace*{1.6em}(ii) Is he right? Explain. \;
Start: \eng{``Mr Bowen is \ldots\ because \ldots''}
\fr{\hspace*{1.6em}A-t-il raison~? Explique, en anglais.}
\qlines{2}

% -- A3 ---------------------------------------------------------------------
% Two lists, two directions. In the long jump the winner is the BIGGEST number
% and it is the shortest-looking one (4.8). In the swim the winner is the
% SMALLEST number and the loser is the shortest-looking one (31.4). (e) asks
% for FOURTH, not the deck's "third": with five lanes the third is the middle
% from either end and cannot catch a student counting from the wrong end.
\qhead{A3}{Who won?\frin{Qui a gagn\'e~?}}

\begin{wordhelp}[Word help --- the words that tell you which end]
  \begin{tabularx}{\linewidth}{@{}l@{\hspace{16pt}}X@{}}
    \eng{longest} & the biggest distance \frin{le plus long} \\
    \eng{fastest} & the \textbf{smallest} time. Fast runners and swimmers take
      \textbf{less} time \frin{le plus rapide = le plus petit temps} \\
    \eng{fourth} & number 4 in the order. \emph{Fourth fastest}: the fourth smallest
      time \frin{quatri\`eme} \\
  \end{tabularx}
\end{wordhelp}

\vspace{4pt}
\renewcommand{\arraystretch}{1.25}%
\noindent
\begin{minipage}[t]{0.47\textwidth}
  \centering\textbf{Long jump}\par\vspace{3pt}
  \begin{tabular}{|C{2.0cm}|C{2.6cm}|}
    \hline
    \textbf{Jumper} & \textbf{Distance (m)} \\ \hline
    A & 4.8  \\ \hline
    B & 4.08 \\ \hline
    C & 4.75 \\ \hline
    D & 4.79 \\ \hline
    E & 4.5  \\ \hline
  \end{tabular}
\end{minipage}\hfill
\begin{minipage}[t]{0.47\textwidth}
  \centering\textbf{50 m swim}\par\vspace{3pt}
  \begin{tabular}{|C{2.0cm}|C{2.6cm}|}
    \hline
    \textbf{Lane} & \textbf{Time (s)} \\ \hline
    1 & 31.4  \\ \hline
    2 & 30.95 \\ \hline
    3 & 31.09 \\ \hline
    4 & 30.9  \\ \hline
    5 & 31.12 \\ \hline
  \end{tabular}
\end{minipage}
\renewcommand{\arraystretch}{1.45}%

\vspace{8pt}
\begin{enumerate}[label=(\alph*), itemsep=8pt]
  \item Write the five jumps in order of size, starting with the smallest.
        \fr{Range les cinq sauts par ordre de taille, du plus petit au plus grand.}
        \ansline[9.0cm]
  \item The \textbf{longest} jump wins. Which jumper won? \hfill \blank[2.6cm]
        \fr{Le saut le plus long gagne. Qui a gagn\'e~?}
  \item Write the five swim times in order of size, starting with the smallest.
        \fr{Range les cinq temps de natation, du plus petit au plus grand.}
        \ansline[9.0cm]
  \item Which lane was the \textbf{fastest}? \frin{le plus rapide} \hfill
        \blank[2.6cm]
  \item Which lane came \textbf{fourth}? \frin{quatri\`eme} \hfill \blank[2.6cm]
  \item Which lane was the \textbf{slowest}? \frin{le plus lent} \hfill
        \blank[2.6cm]
\end{enumerate}

% ===========================================================================
% SECTION B -- MATHEMATICS 4.2
% ===========================================================================
\hwsection{Section B \textperiodcentered\ Mathematics 4.2 --- Adding and Subtracting}

% -- B1 ---------------------------------------------------------------------
\qhead{B1}{Use a written method}

\begin{remember}[Line up the points]
  A \key{written method} is in columns on paper. Write the numbers with the
  \key{points in a line}. Fill the gaps with \key{zeros}. Work from the right, like
  whole numbers. The point in the answer goes \key{straight down}.
  \fr{M\'ethode \'ecrite (\eng{written method})~: en colonnes. Aligne les points.
  Remplis les trous avec des z\'eros. Le point de la r\'eponse descend tout droit
  (\eng{straight down}).}
\end{remember}

\vspace{4pt}
\noindent
\wsq[2.4cm]{a}{5.36 + 2.7}\hfill\wsq[2.4cm]{b}{14.8 + 3.65}\hfill\wsq[2.4cm]{c}{0.45 + 6 + 1.9}

\vspace{4pt}
\noindent
\wsq[2.4cm]{d}{9.4 - 3.75}\hfill\wsq[2.4cm]{e}{12 - 4.36}\hfill\wsq[2.4cm]{f}{1 - 0.65}

% -- B2 ---------------------------------------------------------------------
% The 4.2 mental move: round to a whole number, then give back. The give-back
% sign is the whole question -- (b) and (d) give back with a PLUS.
\qhead{B2}{Use a mental method}

\noindent A \eng{mental method} is in your head. Change one number to a whole
number, then give back the difference. Write the easier calculation, then the answer.
The first one is done for you.
\fr{Calcul mental (\eng{mental method})~: dans ta t\^ete. Change un nombre en nombre
entier, puis rends la diff\'erence. \'Ecris le calcul plus facile, puis la r\'eponse.
Le premier est fait pour toi.}

\vspace{4pt}
\noindent
\begin{tabularx}{\textwidth}{|C{3.0cm}|L|C{2.6cm}|}
  \hline
  \textbf{Calculation} & \textbf{The easier way} & \textbf{Answer} \\
  \hline
  $5.8 + 2.9$ \tall & $5.8 + 3 - 0.1$ & $8.7$ \\ \hline
  $3.9 + 2.5$ \tall & & \\ \hline
  $7.3 - 1.9$ \tall & & \\ \hline
  $6.99 + 3.45$ \tall & & \\ \hline
  $12.6 - 4.98$ \tall & & \\ \hline
\end{tabularx}

% ===========================================================================
% SECTION C -- MATHEMATICS 4.3 AND 4.4
% ===========================================================================
\hwsection{Section C \textperiodcentered\ Mathematics 4.3--4.4 --- Multiplying and Dividing}

% -- C1 ---------------------------------------------------------------------
\qhead{C1}{Point out, point back}

\begin{remember}[Multiplying a decimal]
  Point out: $32 \times 7 = 224$. \; Point back, \key{same number of decimal
  places}: $3.2 \times 7 = 22.4$ \\
  \key{Estimate} to check: \; $3 \times 7 = 21$, and $22.4$ is close to 21.
  \fr{Enl\`eve le point et multiplie. Remets le point avec le m\^eme nombre de
  d\'ecimales (\eng{decimal places}). Estime (\eng{estimate}) pour v\'erifier.}
\end{remember}

\vspace{2pt}
\noindent Fill in the table. The first one is done for you.
\fr{Compl\`ete le tableau. Le premier est fait pour toi.}

\vspace{4pt}
\noindent
\begin{tabularx}{\textwidth}{|C{2.6cm}|L|C{3.5cm}|C{2.4cm}|}
  \hline
  \textbf{Calculation} & \textbf{Estimate} & \textbf{Without the point} &
  \textbf{Answer} \\
  \hline
  $4.6 \times 3$ \tall & $5 \times 3 = 15$ & $46 \times 3 = 138$ & $13.8$ \\ \hline
  $6.4 \times 5$ \tall & & & \\ \hline
  $0.7 \times 8$ \tall & & & \\ \hline
  $2.35 \times 4$ \tall & & & \\ \hline
  $8.15 \times 2$ \tall & & & \\ \hline
\end{tabularx}

% -- C2 ---------------------------------------------------------------------
\qhead{C2}{The point goes straight up}

\begin{remember}[Dividing a decimal]
  Use short division. The point in the answer goes \key{straight up}. \\
  If a digit is too small to divide, write \key{0} and carry it. \\
  Check with the \key{inverse calculation}: \; $8.24 \div 4 = 2.06$ \; because \;
  $2.06 \times 4 = 8.24$.
  \fr{Le point monte tout droit (\eng{straight up}). Si un chiffre est trop petit pour
  \^etre divis\'e, \'ecris 0. V\'erifie avec l'op\'eration inverse (\eng{inverse
  calculation}).}
\end{remember}

\vspace{4pt}
\noindent
\wsq[2.45cm]{a}{7.56 \div 3}\hfill\wsq[2.45cm]{b}{6.24 \div 6}\hfill\wsq[2.45cm]{c}{0.84 \div 7}

\vspace{4pt}
\noindent
\wsq[2.45cm]{d}{5.4 \div 4}\hfill\wsq[2.45cm]{e}{16.2 \div 5}\hfill\wsq[2.45cm]{f}{8.08 \div 8}

\vspace{4pt}
% One line, so (g) cannot be split from its blanks (it was, on the first build).
\noindent \textbf{(g)} Check \textbf{(b)} with the inverse calculation. \hfill
\blank[2.2cm] $\times$ \blank[1.4cm] $=$ \blank[2.2cm]
\fr{(g) V\'erifie (b) avec l'op\'eration inverse.}

% -- C3 ---------------------------------------------------------------------
\qhead{C3}{One fact, five answers}

\noindent You are given one fact. \textbf{Do not multiply again.} Use the fact, and
an estimate to decide where the point goes.
\fr{Tu as un seul calcul. \textbf{Ne multiplie pas encore une fois.} Utilise-le, avec
une estimation, pour savoir o\`u va le point.}

\vspace{6pt}
\begin{center}
\expr{38 \times 27 = 1026}
\end{center}

\vspace{4pt}
\begin{multicols}{3}
\begin{enumerate}[label=(\alph*), itemsep=12pt]
  \item $3.8 \times 27 = \blank[1.7cm]$
  \item $0.38 \times 27 = \blank[1.7cm]$
  \item $38 \times 2.7 = \blank[1.7cm]$
  \item $102.6 \div 27 = \blank[1.7cm]$
  \item $10.26 \div 27 = \blank[1.7cm]$
\end{enumerate}
\end{multicols}

\vspace{2pt}
% A minipage: on one build the two lines landed on two different pages.
\noindent\begin{minipage}{\textwidth}
\textbf{(f)} Two of your answers are the same. Which two, and why?
\fr{Deux de tes r\'eponses sont pareilles. Lesquelles, et pourquoi~?}
\wlines{2}
\end{minipage}

% ===========================================================================
% SECTION D -- MATHEMATICS 4.5
% ===========================================================================
\hwsection{Section D \textperiodcentered\ Mathematics 4.5 --- Making Calculations Easier}

% -- D1 ---------------------------------------------------------------------
\qhead{D1}{Decimals as fractions}

\begin{wordhelp}[Word help]
  \begin{tabularx}{\linewidth}{@{}l@{\hspace{16pt}}X@{}}
    \eng{equivalent} & equal in value, but written differently. $0.3$ and
      $\frac{3}{10}$ are equivalent \frin{\'equivalent} \\
    \eng{tenths} \sep \eng{hundredths} & the column of the \textbf{last} digit is the
      bottom number \frin{dixi\`emes, centi\`emes} \\
  \end{tabularx}
\end{wordhelp}

\vspace{4pt}
\noindent Write each decimal as a fraction, and each fraction as a decimal.
\fr{\'Ecris chaque d\'ecimal en fraction, et chaque fraction en d\'ecimal.}

\vspace{4pt}
\begin{multicols}{3}
\begin{enumerate}[label=(\alph*), itemsep=14pt]
  \item $0.7 = \blank[1.6cm]$
  \item $0.03 = \blank[1.6cm]$
  \item $0.31 = \blank[1.6cm]$
  \item $0.6 = \blank[1.6cm]$
  \item $\frac{8}{100} = \blank[1.6cm]$
  \item $\frac{51}{100} = \blank[1.6cm]$
\end{enumerate}
\end{multicols}

% -- D2 ---------------------------------------------------------------------
% Calculation -> English in rows 2-3, English -> calculation in rows 4-6. The
% English point is that "of" means x.
\qhead{D2}{Half of, a quarter of, a tenth of}

\begin{remember}[In English, ``of'' often means $\times$]
  $\times\, 0.5$ = \key{half of} = $\div\, 2$ \qquad
  $\times\, 0.25$ = \key{a quarter of} = $\div\, 4$ \qquad
  $\times\, 0.1$ = \key{a tenth of} = $\div\, 10$
  \fr{\eng{half of} = la moiti\'e de \quad \eng{a quarter of} = le quart de \quad
  \eng{a tenth of} = le dixi\`eme de}
\end{remember}

\vspace{2pt}
\noindent Fill in the table. The first one is done for you.
\fr{Compl\`ete le tableau. Le premier est fait pour toi. \'Ecris la colonne
\fq{In English} en anglais.}

\vspace{4pt}
\noindent
\begin{tabularx}{\textwidth}{|C{3.4cm}|L|C{2.6cm}|}
  \hline
  \textbf{Calculation} & \textbf{In English} & \textbf{Answer} \\
  \hline
  $48 \times 0.5$ \tall & half of 48 & $24$ \\ \hline
  $36 \times 0.25$ \tall & & \\ \hline
  $70 \times 0.1$ \tall & & \\ \hline
  \tall & a quarter of 2.8 & \\ \hline
  \tall & a tenth of 4.5 & \\ \hline
\end{tabularx}

% -- D3 ---------------------------------------------------------------------
\qhead{D3}{Multiplying by a decimal}

\begin{remember}[Divide first, then multiply]
  $\times\, 0.2$ means $\div\, 10$, then $\times\, 2$. \quad
  $\times\, 0.03$ means $\div\, 100$, then $\times\, 3$. \\
  So \; $250 \times 0.2 = 25 \times 2 = 50$. \;
  Multiplying by a decimal \key{less than 1} makes a number \key{smaller}.
  \fr{Divise d'abord, puis multiplie. Multiplier par un d\'ecimal plus petit que 1
  rend le nombre plus petit (\eng{smaller}).}
\end{remember}

\vspace{4pt}
\begin{multicols}{2}
\begin{enumerate}[label=(\alph*), itemsep=12pt]
  \item $40 \times 0.3 = \blank[1.4cm] \times \blank[1.0cm] = \blank[1.4cm]$
  \item $150 \times 0.2 = \blank[1.4cm] \times \blank[1.0cm] = \blank[1.4cm]$
  \item $600 \times 0.04 = \blank[1.4cm] \times \blank[1.0cm] = \blank[1.4cm]$
  \item $90 \times 0.7 = \blank[1.4cm] \times \blank[1.0cm] = \blank[1.4cm]$
  \item $1200 \times 0.05 = \blank[1.4cm] \times \blank[1.0cm] = \blank[1.4cm]$
  \item $80 \times 0.6 = \blank[1.4cm] \times \blank[1.0cm] = \blank[1.4cm]$
\end{enumerate}
\end{multicols}

% -- D4 ---------------------------------------------------------------------
\qhead{D4}{Times 9, times 11, and partitioning}

\begin{remember}[Three quick ways]
  $\times\, 9$: do $\times\, 10$, then \key{take away} one lot. \;
  $3.8 \times 9 = 38 - 3.8 = 34.2$ \\
  $\times\, 11$: do $\times\, 10$, then \key{add} one lot. \;
  $3.8 \times 11 = 38 + 3.8 = 41.8$ \\
  \key{Partitioning}: split the number. \;
  $4.3 \times 6 = 4 \times 6 + 0.3 \times 6 = 24 + 1.8 = 25.8$
  \fr{$\times\, 9$~: fais $\times\, 10$, puis enl\`eve le nombre une fois. $\times\, 11$~:
  fais $\times\, 10$, puis ajoute-le une fois. \eng{Partitioning} = d\'ecomposer le
  nombre.}
\end{remember}

\vspace{2pt}
\noindent Write the quick way, then the answer.
\fr{\'Ecris la m\'ethode rapide, puis la r\'eponse.}

\vspace{4pt}
\begin{enumerate}[label=(\alph*), itemsep=10pt]
  \item $4.6 \times 9 = $ \blank[8.4cm] $=$ \blank[2.2cm]
  \item $2.3 \times 11 = $ \blank[8.2cm] $=$ \blank[2.2cm]
  \item $6.2 \times 7 = $ \blank[8.4cm] $=$ \blank[2.2cm]
  \item $3.8 \times 5 = $ \blank[8.4cm] $=$ \blank[2.2cm]
\end{enumerate}

% -- D5 ---------------------------------------------------------------------
% One trap from each lesson, which makes this a five-minute diagnosis of the
% whole unit. See the key.
\qhead{D5}{Mr Bowen's homework}

\noindent Mr Bowen gave himself 4 out of 4. Only \textbf{one} of these is right. Find
it, and correct the other three.
\fr{Mr Bowen s'est donn\'e 4 sur 4. Une seule r\'eponse est juste. Trouve-la, et
corrige les trois autres.}

\vspace{4pt}
\begin{multicols}{2}
\begin{enumerate}[label=(\alph*), itemsep=12pt]
  \item $0.4 \times 7 = 0.28$ \; \blank[2.2cm]
  \item $5.09 > 5.1$ \; \blank[2.2cm]
  \item $7.2 \div 6 = 1.2$ \; \blank[2.2cm]
  \item $50 \times 0.3 = 150$ \; \blank[2.2cm]
\end{enumerate}
\end{multicols}

% ===========================================================================
% SECTION E -- THE UNIT 4 REVIEW
% ===========================================================================
\hwsection{Section E \textperiodcentered\ Unit 4 Review --- Bigger Numbers, Real Problems}

% -- E1 ---------------------------------------------------------------------
\qhead{E1}{Two-digit numbers}

\begin{remember}[Estimate, or write the table]
  $\times$ a 2-digit number: take the point out, multiply, then use an
  \key{estimate} to put it back. \; $4.7 \times 32$: \; $47 \times 32 = 1504$, and
  $5 \times 30 = 150$, so \key{150.4}. \\
  $\div$ a 2-digit number: write its \key{times table} first, then use short
  division. The point goes straight up.
  \fr{$\times$ un nombre \`a 2 chiffres~: l'estimation (\eng{estimate}) place le point.
  $\div$ un nombre \`a 2 chiffres~: \'ecris d'abord sa table de multiplication
  (\eng{times table}).}
\end{remember}

\vspace{4pt}
\noindent
\wsh{a}{5.8 \times 15}\hfill\wsh{b}{41.6 \div 13}

% -- E2 ---------------------------------------------------------------------
% THE HEART OF THE MATHS. The numbers are easy on purpose; every item is
% decided by one English word. (c) is "Mr Bowen and two friends" = THREE, the
% deck's slide 20 trap with new numbers. (e) is "subtract X from Y", which
% puts the numbers in the opposite order to the sentence.
\qhead{E2}{From English to a calculation}

\begin{wordhelp}[Word help --- the words that choose the operation]
  \begin{tabularx}{\linewidth}{@{}l@{\hspace{16pt}}X@{}}
    \eng{change} & in a shop, the money you get \textbf{back} when you pay too much.
      Take away \frin{la monnaie} \\
    \eng{share equally} & give everyone the same amount. Divide --- and count the
      people carefully \frin{partager en parts \'egales} \\
    \eng{how much heavier} & the difference. Take away \frin{combien plus lourd} \\
    \eng{subtract 2 from 9} & $9 - 2$. The \textbf{second} number comes
      \textbf{first} \frin{soustrais 2 de 9} \\
  \end{tabularx}
\end{wordhelp}

\vspace{2pt}
\noindent Write the \textbf{calculation first}, then the answer. Give units.
\fr{\'Ecris d'abord le calcul (\eng{calculation}), puis la r\'eponse. Donne les
unit\'es.}

\begin{enumerate}[label=(\alph*), itemsep=6pt]
  \item Mr Bowen pays for a 3.65-dollar sandwich with a 5-dollar note. How much
        change does he get?
        \fr{Mr Bowen paie un sandwich \`a 3.65 dollars avec un billet de 5 dollars.
        Combien de monnaie (\eng{change}) re\c{c}oit-il~?}
        \calcans
  \item A mango costs 1.20 dollars. How much do 6 mangoes cost?
        \fr{Une mangue co\^ute 1.20 dollar. Combien co\^utent 6 mangues~?}
        \calcans
  \item Mr Bowen and two friends share a 13.50-dollar pizza equally. How much does
        each person pay?
        \fr{Mr Bowen et deux amis partagent une pizza \`a 13.50 dollars en parts
        \'egales (\eng{share equally}). Combien paie chaque personne~?}
        \calcans
  \item A bag of rice weighs 2.4 kg. A bag of sugar weighs 0.85 kg. How much
        heavier is the rice?
        \fr{Un sac de riz p\`ese 2.4 kg. Un sac de sucre p\`ese 0.85 kg. Combien le
        riz est-il plus lourd (\eng{how much heavier})~?}
        \calcans
  \item Subtract 2.6 from 9.1.
        \fr{Soustrais 2.6 de 9.1 (\eng{subtract 2.6 from 9.1}).}
        \calcans
\end{enumerate}

% -- E3 ---------------------------------------------------------------------
\qhead{E3}{Money has two decimal places}

\begin{remember}[Money]
  Money is written to \key{2 decimal places}. \;
  4.5 dollars is written \key{4.50} dollars. \\
  $10 \div 3 = 3.333\ldots$, so three friends get \key{3.33} dollars each, and
  1 cent is left over.
  \fr{L'argent s'\'ecrit avec 2 d\'ecimales (\eng{2 decimal places})~: 4.5 dollars
  s'\'ecrit 4.50 dollars.}
\end{remember}

\vspace{6pt}
\noindent Write each one as money.
\fr{\'Ecris chaque montant comme de l'argent.}

\vspace{6pt}
\noindent\hspace*{1.6em}\textbf{(a)} 4.8 dollars \; \blank[2.0cm] \hfill
\textbf{(b)} 9.125 dollars \; \blank[2.0cm] \hfill
\textbf{(c)} $15 \div 4$ dollars \; \blank[2.0cm]
\par\vspace{10pt}

\vspace{2pt}
\noindent \textbf{(d)} Three friends share 25 dollars equally. How much does each
friend get? How many cents are left over?
\fr{(d) Trois amis partagent 25 dollars en parts \'egales. Combien re\c{c}oit chacun~?
Combien de cents reste-t-il (\eng{left over})~?}

\vspace{5pt}
\noindent\hspace*{1.6em}Each friend: \blank[3.0cm] dollars \hfill
Left over: \blank[2.4cm] cents

% -- E4 ---------------------------------------------------------------------
% (b) is the deadpan one. It needs nothing untaught: "a quarter of" is / 4
% (D2), and 0.500 / 4 needs the zeros written (C2).
\qhead{E4}{Word problems}

\begin{enumerate}[label=(\alph*), itemsep=6pt]
% Each item is a minipage, so its question cannot be left on one page with
% its box on the next (it was, in the French copy).
  \item \begin{minipage}[t]{\linewidth}
        Mr Bowen buys 3 notebooks at 2.45 dollars each and one pen for 1.80
        dollars. He pays with a 10-dollar note. How much change does he get?
        \fr{Mr Bowen ach\`ete 3 cahiers \`a 2.45 dollars chacun et un stylo \`a 1.80
        dollar. Il paie avec un billet de 10 dollars. Combien de monnaie
        re\c{c}oit-il~?}
        \begin{workspace}[2.7cm]\end{workspace}
        \end{minipage}
  \item \begin{minipage}[t]{\linewidth}
        Mr Bowen buys 0.5 kg of grapes. He eats a quarter of them on the walk home,
        which he calls ``quality control''. How many kilograms of grapes arrive at
        his house?
        \fr{Mr Bowen ach\`ete 0.5 kg de raisin. Il en mange un quart (\eng{a
        quarter}) en rentrant, pour le \fq{contr\^ole qualit\'e}. Combien de kilos de
        raisin arrivent chez lui~?}
        \begin{workspace}[2.7cm]\end{workspace}
        \end{minipage}
\end{enumerate}

% -- E5 ---------------------------------------------------------------------
\qhead{E5}{Now you write the question}

\begin{yourturn}[The rules]
  A \textbf{real} situation --- a shop, a recipe, a race, a bag on a scale. Full
  sentences, ending in a \textbf{question}. Use one of the \eng{words} from E2.
  Give \textbf{units}. Your question must come out at the answer given.
  \fr{Une situation r\'eelle, des phrases compl\`etes, une question \`a la fin, un mot
  de E2, des unit\'es. \textbf{\'Ecris ton probl\`eme en anglais.}}
\end{yourturn}

\vspace{4pt}
\noindent Write a word problem for \quad \expr{12.60 \div 4 = 3.15}
\fr{\'Ecris un probl\`eme (\eng{word problem}) pour ce calcul, en anglais.}
\qlines{3}

% ===========================================================================
% SECTION F -- SCIENCE 3.1 IN DECIMALS
% ===========================================================================
\hwsection[cbgreen]{Section F \textperiodcentered\ Science 3.1 in decimals --- Mr Bowen in Space}

\begin{remember}[Weight and mass, from Science 3.1]
  \key{weight (N) = mass (kg) $\times$ strength of gravity (N/kg)} \\
  On Earth, the strength of gravity is \key{10 N/kg}. Say it: \eng{``ten newtons
  per kilogram''}. \\
  \key{Mass does not change} from planet to planet. Weight does.
  \fr{Poids (\eng{weight}, en N) = masse (\eng{mass}, en kg) $\times$ intensit\'e de la
  gravit\'e (\eng{strength of gravity}). La masse ne change pas d'une plan\`ete \`a
  l'autre~; le poids, oui.}
\end{remember}

% -- F1 ---------------------------------------------------------------------
\qhead[cbgreen]{F1}{How strong is gravity?}

% A minipage, so the heading and its table cannot be split (they were).
\noindent\begin{minipage}{\textwidth}
These are the strengths of gravity from the lesson. A bigger number pulls harder.
\fr{Voici l'intensit\'e de la gravit\'e de la le\c{c}on. Un nombre plus grand tire
plus fort.}

\vspace{8pt}
\centering
\setlength{\tabcolsep}{5pt}%
\begin{tabular}{|l|*{7}{C{1.70cm}|}}
  \hline
  \textbf{Place} & Moon & Mars & Venus & Uranus & Earth & Neptune & Jupiter \\
  \hline
  \textbf{N/kg} & \textbf{1.6} & \textbf{3.7} & \textbf{8.9} & \textbf{8.7} &
  \textbf{10} & \textbf{11} & \textbf{25} \\
  \hline
\end{tabular}
\end{minipage}

\vspace{4pt}
\begin{enumerate}[label=(\alph*), itemsep=8pt]
  \item Write the seven strengths in order of size, starting with the smallest.
        \fr{Range les sept valeurs, de la plus petite \`a la plus grande.}
        \ansline[11.0cm]
  \item Which pulls harder, \textbf{Venus} or \textbf{Uranus}? By how much?
        \hfill \blank[5.4cm]
        \fr{Laquelle tire plus fort, V\'enus ou Uranus~? De combien~?}
  \item How much stronger is gravity on the \textbf{Earth} than on the
        \textbf{Moon}? \hfill \blank[4.0cm]
        \fr{De combien la gravit\'e est-elle plus forte sur la Terre que sur la
        Lune~?}
\end{enumerate}

% -- F2 ---------------------------------------------------------------------
\qhead[cbgreen]{F2}{Mr Bowen's space backpack}

\noindent Mr Bowen's space backpack has a mass of \textbf{15 kg}. \textbf{(a)} Find its weight in
each place. Estimate first. The first one is done for you.
\fr{Le sac \`a dos spatial de Mr Bowen a une masse de 15 kg. (a) Trouve son poids
(\eng{weight}) \`a chaque endroit. Estime d'abord.}

\vspace{4pt}
\noindent
\begin{tabularx}{\textwidth}{|C{2.4cm}|C{3.0cm}|L|C{2.8cm}|}
  \hline
  \textbf{Place} & \textbf{Gravity (N/kg)} & \textbf{Calculation} &
  \textbf{Weight (N)} \\
  \hline
  Earth \tall   & 10  & $15 \times 10$ & $150$ \\ \hline
  Moon \tall    & 1.6 & & \\ \hline
  Mars \tall    & 3.7 & & \\ \hline
  Venus \tall   & 8.9 & & \\ \hline
  Jupiter \tall & 25  & & \\ \hline
\end{tabularx}

\vspace{6pt}
\noindent \textbf{(b)} How much \textbf{less} does the backpack weigh on the Moon than
on the Earth? \hfill \blank[2.6cm]
\fr{(b) De combien le sac p\`ese-t-il moins (\eng{less}) sur la Lune que sur la
Terre~?}

% -- F3 ---------------------------------------------------------------------
% THE KEYSTONE. The light ball has the fastest single time (0.39 s), so a
% student who looks at one number says the light ball is faster. Both means
% are 1.80 / 4 = 0.45 s. The decimals answer the science question.
\qhead[cbgreen]{F3}{Do heavy things fall faster?}

\noindent Mr Bowen drops a \textbf{heavy} ball and a \textbf{light} ball, the same
size, from 1 metre. He times each fall with a stopwatch, four times.
\fr{Mr Bowen l\^ache une balle lourde (\eng{heavy}) et une balle l\'eg\`ere
(\eng{light}) de m\^eme taille, d'un m\`etre. Il chronom\`etre chaque chute quatre
fois.}

\vspace{4pt}
\begin{center}
\begin{tabular}{|l|C{1.6cm}|C{1.6cm}|C{1.6cm}|C{1.6cm}|}
  \hline
  & \textbf{Drop 1} & \textbf{Drop 2} & \textbf{Drop 3} & \textbf{Drop 4} \\
  \hline
  \textbf{Heavy ball} (s) & 0.48 & 0.41 & 0.50 & 0.41 \\ \hline
  \textbf{Light ball} (s) & 0.44 & 0.52 & 0.39 & 0.45 \\ \hline
\end{tabular}
\end{center}

\begin{wordhelp}[Word help]
  \begin{tabularx}{\linewidth}{@{}l@{\hspace{16pt}}X@{}}
    \eng{conjecture} & a guess that you can test \frin{une conjecture} \\
    \eng{mean} & add up all the values, then divide by how many there are
      \frin{la moyenne} \\
    \eng{reaction time} & how long you take to move. It is about \textbf{0.2 s}
      \frin{le temps de r\'eaction} \\
  \end{tabularx}
\end{wordhelp}

\vspace{4pt}
\begin{enumerate}[label=(\alph*), itemsep=8pt]
  \item Find the \textbf{fastest} time in the table. Which ball was it?
        \hfill \blank[4.0cm]
        \fr{Trouve le temps le plus rapide (\eng{fastest}). C'est quelle balle~?}
  \item Make a \textbf{conjecture} from (a): which ball falls faster?
        \hfill \blank[4.0cm]
        \fr{Fais une conjecture avec (a)~: quelle balle tombe plus vite~?}
  \item Now test it. Find the \textbf{mean} time for each ball.
        \fr{Teste-la. Trouve le temps moyen (\eng{mean}) de chaque balle.}
        \par\vspace{4pt}
        \begin{minipage}[t]{0.47\linewidth}
          \textbf{Heavy ball}
          \begin{workspace}[3.3cm]\end{workspace}
        \end{minipage}\hfill
        \begin{minipage}[t]{0.47\linewidth}
          \textbf{Light ball}
          \begin{workspace}[3.3cm]\end{workspace}
        \end{minipage}
  \item Do heavy things fall faster? Use your two means. Answer in a full sentence.
        \fr{Les objets lourds tombent-ils plus vite~? Utilise tes deux moyennes.
        R\'eponds par une phrase compl\`ete, en anglais.}
        \qlines{3}
  \item The four times for the \textbf{same} ball are all different. Why? Use the
        words \eng{reaction time}.
        \fr{Les quatre temps de la m\^eme balle sont tous diff\'erents. Pourquoi~?
        Utilise les mots \eng{reaction time}, en anglais.}
        \qlines{3}
\end{enumerate}

% -- F4 ---------------------------------------------------------------------
% The deck's "Mr Bowen in Space" run backwards: weight given, find the
% strength of gravity, then read the place off F1. 296 / 80 = 29.6 / 8 = 3.7,
% or try each place: 80 x 3.7 = 296. Mars is the only 3.7 in F1 on purpose.
\qhead[cbgreen]{F4}{Mr Bowen moves house}

\noindent Mr Bowen's mass is \textbf{80 kg}. He would like to weigh exactly
\textbf{296 N} without eating any less. Which place in \textbf{F1} should he move
to? Show your calculation.
\fr{Mr Bowen a une masse de 80 kg. Il veut peser exactement 296 N sans manger moins.
O\`u doit-il d\'em\'enager (choisis dans F1)~? Montre ton calcul.}

\begin{workspace}[4.0cm]\end{workspace}

\noindent He has started packing. What will his \textbf{mass} be when he gets there?
\hfill \blank[3.0cm]
\fr{Il a commenc\'e ses valises. Quelle sera sa masse (\eng{mass}) l\`a-bas~?}

% ===========================================================================
% ANSWER KEY -- TEACHER COPY ONLY
% ===========================================================================
\ifkey
\clearpage

\hwsection[cbred]{Answer Key --- Teacher Copy}

\linespread{1.0}\selectfont
\setlist[enumerate]{itemsep=2pt, topsep=3pt, leftmargin=*}
\setlist[itemize]{itemsep=2pt, topsep=3pt, leftmargin=*}

\qhead[cbred]{A1}{Put each number on the line}
\noindent A on the 2nd small line \sep B on the 9th \sep C halfway between the 4th
and 5th \sep D halfway between 5 and the 1st. \\
\textbf{(b)} $5.05, \; 5.2, \; 5.45, \; 5.9$

\smallskip
\noindent{\small D is the one to watch. A student who puts $5.05$ \emph{after} $5.2$
has read it as $5.5$ --- the zero after the point was not said --- and that is the
whole of 4.1 in one mark. The line is there so they can see $5.05$ is the first thing
after 5. If (b) is right but the line is wrong, they ordered by the rule and did not
trust the picture; that is fine.}

\qhead[cbred]{A2}{Which is bigger?}
\noindent (a) $>$ \quad (b) $>$ \quad (c) $<$ \quad (d) $>$ \quad (e) $=$ \quad
(f) $<$ \\
\textbf{(g)(i)} six point one eight \sep six point two. \\
\textbf{(ii)} He is wrong. $6.2 = 6.20$, and $6.20 > 6.18$: two tenths is more than
one tenth.

\smallskip
\noindent{\small (g) is the deck's slide 3 with new numbers, and (i) is the point of
it. Anyone who writes ``six point eighteen'' has found the cause of the mistake by
themselves: English reads the digits after the point one at a time, and
``eighteen'' is the French \emph{virgule dix-huit} (and the Vietnamese) leaking
through. Read (i) before marking (ii). \\
(d) is the one that goes wrong the other way: $9.99$ has more digits and more nines,
and $12.3$ is bigger because the whole-number part is checked first. \\
(f) needs \emph{two} zeros of filling: $4.125$ against $4.130$.}

\qhead[cbred]{A3}{Who won?}
\noindent (a) $4.08, \; 4.5, \; 4.75, \; 4.79, \; 4.8$ \quad (b) \textbf{A} \\
(c) $30.9, \; 30.95, \; 31.09, \; 31.12, \; 31.4$ \quad (d) \textbf{Lane 4} \quad
(e) \textbf{Lane 5} \quad (f) \textbf{Lane 1}

\smallskip
\noindent{\small Every answer here is the number that \emph{looks} wrong. The winning
jump, $4.8$, is the shortest-looking number in the list; the slowest swim, $31.4$, is
the shortest-looking time. A student who orders by length gets both upside down. \\
(e) asks for \emph{fourth}, not the deck's \emph{third}, on purpose: with five lanes
the third is the middle one from either end, so it cannot catch anyone counting from
the wrong end. Fourth can. Expect Lane 2, which is fourth counted from the
\emph{slowest}. (d) is the deck's ``fastest means smallest'' --- ask anyone who wrote
Lane 1 whether a fast swimmer takes more time or less.}

\qhead[cbred]{B1}{Use a written method}
\noindent (a) $8.06$ \quad (b) $18.45$ \quad (c) $8.35$ \quad (d) $5.65$ \quad
(e) $7.64$ \quad (f) $0.35$

\smallskip
\noindent{\small (a) is the alignment check: anyone who lines up the \emph{right-hand
digits} instead of the points adds $536 + 27$ and writes $5.63$. \\
(c) has a whole number with no point at all. $6$ is $6.00$; students who cannot see
where its point is put the 6 under the 5 and get $2.41$. \\
(d), (e) and (f) need the zeros \emph{written} before subtracting: $9.40$, $12.00$,
$1.00$. (f) is the deck's $1 - 0.72$ with new numbers; the usual wrong answer is
$0.45$ (borrowing from the tenths and forgetting the hundredths).}

\qhead[cbred]{B2}{Use a mental method}
\noindent $3.9 + 2.5 = 4 + 2.5 - 0.1 = 6.4$ \\
$7.3 - 1.9 = 7.3 - 2 + 0.1 = 5.4$ \\
$6.99 + 3.45 = 7 + 3.45 - 0.01 = 10.44$ \\
$12.6 - 4.98 = 12.6 - 5 + 0.02 = 7.62$

\smallskip
\noindent{\small The sign of the give-back is the question. When you take away
\emph{too much} you have to give some back, so the subtractions end in $+$. Expect
$5.2$ for the second row and $7.58$ for the fourth --- the right move with the wrong
sign. \\
Row 4 gives back $0.02$, not $0.2$: $5 - 4.98$ is two hundredths. \\
Accept any correct easier way (e.g.\ $3.9 + 2.5 = 3.9 + 2.1 + 0.4$). The column must
show a calculation, not just the answer.}

\qhead[cbred]{C1}{Point out, point back}
\noindent
\begin{tabular}{@{}llll@{}}
  $6.4 \times 5$ & $6 \times 5 = 30$ & $64 \times 5 = 320$ & $\mathbf{32}$ \\
  $0.7 \times 8$ & $1 \times 8 = 8$ & $7 \times 8 = 56$ & $\mathbf{5.6}$ \\
  $2.35 \times 4$ & $2 \times 4 = 8$ & $235 \times 4 = 940$ & $\mathbf{9.4}$ \\
  $8.15 \times 2$ & $8 \times 2 = 16$ & $815 \times 2 = 1630$ & $\mathbf{16.3}$ \\
\end{tabular}

\smallskip
\noindent{\small Rows 1, 3 and 4 end in a zero after the point is put back ($32.0$,
$9.40$, $16.30$). Accept those, and accept the zero dropped: nothing was rounded, as
the deck's own practice answer says ($5.00 = 5$). \\
Row 2 is the deck's vote ($0.3 \times 4$) in disguise. Expect $0.56$ from anyone who
counts \emph{all} the digits instead of the decimal places, and the estimate catches
it: the answer must be near 8. \\
Estimates may differ ($2.35 \times 4$ could be $2.5 \times 4 = 10$); any sensible
rounding is right.}

\qhead[cbred]{C2}{The point goes straight up}
\noindent (a) $2.52$ \quad (b) $\mathbf{1.04}$ \quad (c) $0.12$ \quad
(d) $\mathbf{1.35}$ \quad (e) $\mathbf{3.24}$ \quad (f) $\mathbf{1.01}$ \\
(g) $1.04 \times 6 = 6.24$

\smallskip
\noindent{\small (b) and (f) need a 0 in the answer where a digit does not divide: 6
into 2 will not go. Expect $1.4$ and $1.1$, which is the deck's ``$2.6$ forgot the 0''
exactly. The inverse check in (g) catches it: $1.4 \times 6 = 8.4$, not $6.24$. \\
(d) and (e) run out of digits: write $5.40$ and $16.20$ and keep going. Expect
``$1.3$ remainder 2'' --- the remainder belongs in the hundredths. \\
(c) starts with $0$ in the ones: 7 into 0 will not go, so the answer begins $0.$}

\qhead[cbred]{C3}{One fact, five answers}
\noindent (a) $102.6$ \quad (b) $10.26$ \quad (c) $102.6$ \quad (d) $3.8$ \quad
(e) $0.38$ \\
(f) (a) and (c). Both have one decimal place in total, so both are $1026$ with one
place put back. (Or: $3.8 \times 27$ and $38 \times 2.7$ are both about
$4 \times 27 \approx 100$.)

\smallskip
\noindent{\small Estimates: (a) $4 \times 30 = 120$, so $102.6$ and not $10.26$ or
$1026$. (d) is (a) undone, and (e) is (b) undone --- a student who sees that has the
inverse calculation without being told. \\
Anyone who multiplies $38 \times 27$ again has not read the instruction; it is the
English of ``do not'' and worth one line.}

\qhead[cbred]{D1}{Decimals as fractions}
\noindent (a) $\frac{7}{10}$ \quad (b) $\frac{3}{100}$ \quad (c) $\frac{31}{100}$
\quad (d) $\frac{6}{10}$ (or $\frac{3}{5}$) \quad (e) $0.08$ \quad (f) $0.51$

\smallskip
\noindent{\small (b) and (e) are the same trap in two directions: $0.03 =
\frac{3}{10}$ and $\frac{8}{100} = 0.8$. Both come from reading the digit and
ignoring the column. Send them back to the word-help line: the \emph{last} digit's
column is the bottom number. Simplifying is not required.}

\qhead[cbred]{D2}{Half of, a quarter of, a tenth of}
\noindent a quarter of 36, $9$ \sep a tenth of 70, $7$ \sep $2.8 \times 0.25$,
$0.7$ \sep $4.5 \times 0.1$, $0.45$

\smallskip
\noindent{\small Rows 4 and 5 are the English-to-calculation direction, and the
column is ``calculation'': accept $2.8 \div 4$ as well as $2.8 \times 0.25$, but point
out that the deck wrote it with $\times$ because ``of'' means $\times$. \\
Row 5 is the one to watch: a tenth of $4.5$ is $0.45$, not $0.5$ or $4.4$. \\
Row 4 needs the 0 in the ones: 4 into 2 will not go, so $0.7$ (C2(c) again).}

\qhead[cbred]{D3}{Multiplying by a decimal}
\noindent (a) $4 \times 3 = 12$ \quad (b) $15 \times 2 = 30$ \quad
(c) $6 \times 4 = 24$ \quad (d) $9 \times 7 = 63$ \quad (e) $12 \times 5 = 60$ \quad
(f) $8 \times 6 = 48$

\smallskip
\noindent{\small (c) and (e) divide by \textbf{100}, not 10. Expect $60 \times 4 = 240$
and $120 \times 5 = 600$: the student divided by 10 for a hundredth. \\
Every answer is smaller than the number it started from. Until this unit every
multiplication the class has done made the number bigger (the deck's slide 6); it is
worth asking aloud whether anyone noticed. D5(d) tests it.}

\qhead[cbred]{D4}{Times 9, times 11, and partitioning}
\noindent (a) $46 - 4.6 = 41.4$ \quad (b) $23 + 2.3 = 25.3$ \\
(c) $6 \times 7 + 0.2 \times 7 = 42 + 1.4 = 43.4$ \quad
(d) $3 \times 5 + 0.8 \times 5 = 15 + 4 = 19$

\smallskip
\noindent{\small (a) needs a decimal subtraction with a hidden zero: $46.0 - 4.6$.
Expect $42.4$ from $46 - 4 + 0.6$. \\
(d) is partitioning landing on a whole number: $0.8 \times 5 = 4.0$. Students who
write $0.4$ have counted the decimal places in the question as if they were in the
answer --- $8 \times 5 = 40$, one place back, $4.0$. \\
Accept the column method for (c) and (d), but ask for the partition: the deck taught
it as the faster way.}

\qhead[cbred]{D5}{Mr Bowen's homework}
\noindent (a) \textbf{wrong} --- $2.8$. One decimal place in, one decimal place out;
the estimate is $0.5 \times 7 \approx 3$. \\
(b) \textbf{wrong} --- $5.09 < 5.1$. Fill with zeros: $5.09$ against $5.10$. \\
(c) \textbf{right.} \\
(d) \textbf{wrong} --- $15$. Multiplying by a decimal less than 1 makes it smaller:
$5 \times 3$.

\smallskip
\noindent{\small One trap each from C1, A2, C2 and D3, in that order. Whichever one a
student gets wrong, send them to that question rather than explaining again. \\
(c) is right on purpose, and it is the one with the most room for doubt: $7.2 \div 6$
gives $1.2$ and a student who expects every division in this unit to need a 0 will
``correct'' it to $1.02$.}

\qhead[cbred]{E1}{Two-digit numbers}
\noindent (a) $87$ \quad ($58 \times 15 = 870$; estimate $6 \times 15 = 90$) \\
(b) $3.2$ \quad (13 table: $13 \times 3 = 39$, remainder $2.6 = 13 \times 0.2$)

\smallskip
\noindent{\small (a) lands on a whole number, $87.0$; expect $8.70$ from a student who
put back two places because the question had ``two numbers''. The estimate, 90, is
the fastest correction. \\
For (b) the table is the method --- look for it written down. A student with no table
and the right answer probably knew it; a student with no table and a wrong answer
needs to be told to write the table, not shown the division.}

\qhead[cbred]{E2}{From English to a calculation}
\noindent (a) $5.00 - 3.65 = 1.35$ dollars \quad (b) $1.20 \times 6 = 7.20$ dollars \\
(c) $13.50 \div 3 = 4.50$ dollars \quad (d) $2.4 - 0.85 = 1.55$ kg \\
(e) $9.1 - 2.6 = 6.5$

\smallskip
\noindent{\small \textbf{Mark the calculation, not the answer.} This is the English
question, and the numbers were chosen to be easy. \\
(c) is the deck's slide 20 trap: \emph{Mr Bowen and two friends} is \textbf{three}
people. $13.50 \div 2 = 6.75$ is an English error, not an arithmetic one --- make the
student count the people out loud. \\
(e) puts the numbers in the opposite order to the sentence. A student who writes
$2.6 - 9.1$ has read left to right; ask what ``subtract 2 from 9'' is in the word-help
box. \\
(a) and (b): accept $5 - 3.65$ and $1.2 \times 6$, but the answers are money and
should be $1.35$ and $7.20$ (E3).}

\qhead[cbred]{E3}{Money has two decimal places}
\noindent (a) $4.80$ \quad (b) $9.13$ \quad (c) $3.75$ \\
(d) $25 \div 3 = 8.333\ldots$, so $\mathbf{8.33}$ dollars each, and $25.00 - 24.99 =
0.01$, so \textbf{1 cent} is left over.

\smallskip
\noindent{\small (b) is rounding: the 5 in the thousandths rounds the 2 up. Expect
$9.12$ (cut off, not rounded) and $9.125$ (left alone). \\
(d) is the deck's slide 10 with 25 instead of 10. ``How many cents'' is the English:
some will write $0.01$, which is right in dollars and should be turned into the word
the question used.}

\qhead[cbred]{E4}{Word problems}
\noindent (a) $3 \times 2.45 = 7.35$; \; $7.35 + 1.80 = 9.15$; \; $10.00 - 9.15 =
\mathbf{0.85}$ dollars. \\
(b) A quarter of $0.5$ is $0.5 \div 4 = 0.125$ kg. \; $0.5 - 0.125 = \mathbf{0.375}$
kg.

\smallskip
\noindent{\small (a) is three steps, and the usual slip is to stop after two and give
$9.15$ as ``the change''. Ask what the question asked for. \\
(b) is the deadpan one and the hardest item in E. $0.500 \div 4$ needs two zeros
written to finish, and the question asks what \emph{arrives}, not what he ate. Expect
$0.125$ as the final answer; that is the grapes Mr Bowen ate. Accept $0.5 \times 0.75$
if anyone finds it, but nobody is expected to.}

\qhead[cbred]{E5}{Now you write the question}
\noindent Check, in order: a real situation; full sentences; ends in a question
mark; uses a word from E2; gives units; and the student's own question really comes
out at the answer.

\smallskip
\noindent{\small It wants \emph{share equally} with \textbf{four} people. A good
answer: ``Mr Bowen and three friends share a 12.60-dollar pizza equally. How much does
each person pay?'' \\
Watch for ``Mr Bowen and four friends'', which is five, and is the E2(c) error made by
the student who wrote the question. That one is worth reading aloud to the class. \\
The other common failure is a question whose answer is $50.40$: ``each of 4 friends
has 12.60 dollars --- how much altogether?'' Read it back to them.}

\qhead[cbred]{F1}{How strong is gravity?}
\noindent (a) $1.6, \; 3.7, \; 8.7, \; 8.9, \; 10, \; 11, \; 25$ \\
(b) \textbf{Venus}, by $8.9 - 8.7 = 0.2$ N/kg. \quad
(c) $10 - 1.6 = 8.4$ N/kg.

\smallskip
\noindent{\small (a): $8.7$ and $8.9$ are the pair to check, and $10$ and $11$ after
the two eights. A student ordering by digit count puts $10$, $11$ and $25$ \emph{before}
$8.9$. \\
(c) needs $10.0 - 1.6$; expect $9.6$ or $8.6$ from anyone who did not write the zero.}

\qhead[cbred]{F2}{Mr Bowen's space backpack}
\noindent Moon $15 \times 1.6 = 24$ N \sep Mars $15 \times 3.7 = 55.5$ N \sep
Venus $15 \times 8.9 = 133.5$ N \sep Jupiter $15 \times 25 = 375$ N \\
(b) $150 - 24 = \mathbf{126}$ N

\smallskip
\noindent{\small Estimates: Moon $15 \times 2 = 30$; Mars $15 \times 4 = 60$; Venus
$15 \times 9 = 135$. A student who writes $2.4$ or $240$ for the Moon has the digits
right and the point wrong, and the estimate of 30 settles it. \\
Venus invites $\times 9$: $15 \times 8.9 = 15 \times 9 - 1.5 = 133.5$ is D4 used without
being asked --- worth praising aloud. \\
The mass on every row is still 15 kg; F4's last line asks that, and it is worth
saying here if anyone wrote a different mass for the Moon.}

\qhead[cbred]{F3}{Do heavy things fall faster?}
\noindent (a) $0.39$ s, the \textbf{light} ball. \quad (b) ``The light ball falls
faster.'' \\
(c) Heavy: $0.48 + 0.41 + 0.50 + 0.41 = 1.80$, \; $1.80 \div 4 = \mathbf{0.45}$ s. \;
Light: $0.44 + 0.52 + 0.39 + 0.45 = 1.80$, \; $1.80 \div 4 = \mathbf{0.45}$ s. \\
(d) ``No. Both balls have a mean time of 0.45 s, so they fall at the same speed.'' \\
(e) ``Mr Bowen's reaction time is about 0.2 s, so he starts and stops the stopwatch
a little early or late each time.''

\smallskip
\noindent{\small \textbf{This is the question the packet is built around.} The
fastest single time belongs to the light ball, so the conjecture in (b) is
\emph{supposed} to be wrong. That is what a conjecture is for: the mean tests it and
throws it out. Expect some students to ``fix'' (b) after doing (c); ask them to leave
it and say in (d) that it was wrong. \\
(d) answers the question Science 3.1 opened on (Galileo, the hammer and the
feather), and it lands on the same answer the deck gave --- in decimals, from
numbers. It is worth putting both means on the board next to each other. \\
(e) is slide 43: the times differ by up to $0.11$ s, about half a reaction time. Accept
any sentence that blames the person with the stopwatch and uses ``reaction time''.}

\qhead[cbred]{F4}{Mr Bowen moves house}
\noindent $296 \div 80 = 29.6 \div 8 = 3.7$ N/kg, which is \textbf{Mars}. \;
(Or try each place: $80 \times 3.7 = 296$.) \\
His mass will be \textbf{80 kg} --- mass does not change.

\smallskip
\noindent{\small Either method is fine; trial with F1's list is perfectly good
mathematics and probably the more common route. $296 \div 80$ is the deck's ``Mr
Bowen in Space'' run backwards, and some students will remember 296 N from the lesson
--- that also counts. \\
\textbf{The second line is the science, and it is the one to discuss rather than
mark.} Mass does not change, so it needs no calculation at all, and some students will
distrust that: expect $296$, or $80 \times 3.7$ worked out again. Mercury (also 3.7) is
left out of F1 so that the answer is a single place.}

\fi
\end{document}
